% !PRINTED 545
\weberpage{545}

Wir führen nun die Bezeichnung ein:
\weq{9}{\begin{aligned}
\eta &= \alpha + \alpha_e + \alpha_{2e} + \cdots + \alpha_{(f-1)e}\\
\eta_1 &= \alpha_1 + \alpha_{e+1} + \alpha_{2e+1} + \cdots + \alpha_{(f-1)e+1},\\
&\cdots\cdots\cdots\cdots\cdots\cdots\cdots\cdots\cdots\cdots\\
\eta_{e-1} &= \alpha_{e-1} + \alpha_{2e-1} + \alpha_{3e-1} + \cdots + \alpha_{m-1},
\end{aligned}}

und nennen diese Grössen, wie es \gesperrt{Gauss} in dem speciellen Falle der Kreistheilung gethan hat, die $f$-gliedrigen Perioden der Grössen $\alpha$. Es wird dann
\weq{10}{(\varepsilon,\alpha) = \eta + \varepsilon\eta_1 + \varepsilon^2\eta_2 + \cdots + \varepsilon^{e-1}\eta_{e-1},}
wofür wir auch $(\varepsilon,\eta)$ schreiben können.

Die Anwendung der Permutation $\pi$ auf die Indices von $\alpha$ bringt unter den nach dem Modul $e$ zu nehmenden Indices der Grössen $\eta$ die cyklische Permutation
\weq{11}{\gamma = (0, 1, 2\dots e-1)}
hervor.

Entsprechend der Formel (10) können wir die Formeln (5) und (7) nun auch so schreiben:
\[
(\varepsilon,\eta)^{\nu} = E_0^{(\nu)} + \varepsilon E_1^{(\nu)} + \varepsilon^2 E_2^{(\nu)} + \cdots + \varepsilon^{e-1}E_{e-1}^{(\nu)}
\]
\[
(\varepsilon,\eta)^{\nu}(\varepsilon^{\lambda_1},\eta)^{\nu_1}(\varepsilon^{\lambda_2},\eta)^{\nu_2}\dots = G_0 + \varepsilon G_1 + \varepsilon^2 G_2 + \cdots + \varepsilon^{e-1}G_{e-1},
\]
worin dann
\[
E_k^{(\nu)} = A_k^{(\nu)} + A_{e+k}^{(\nu)} + \cdots + A_{(f-1)e+k}^{(\nu)}
\]
\[
G_k = B_k + B_{e+k} + \cdots + B_{(f-1)e+k}
\]
ist, und $E_k^{(\nu)}$ und $G_k$ ganze homogene Functionen der $\alpha$ sind, die sich auch als Functionen der $\eta$ darstellen lassen.

Die Grössen $\eta_k$, $E_k^{(\nu)}$, $G_k$ bleiben ungeändert, wenn die Permutation $\pi^e$ auf ihre Indices angewandt wird, d. h. wenn der Index um ein Vielfaches von $e$ verändert wird, und wir können also jetzt die Sätze 2. und 3. so vervollständigen:

\gesperrt{4. Ist $\varepsilon$ eine beliebige $e^{\mathrm{te}}$ Einheitswurzel und $e$ ein Theiler von $m$, so erleiden, wenn auf die Indices von $\alpha$ die Permutation $\pi$ ausgeübt wird, die Indices $k$ der Coëfficienten $E_k^{(\nu)}$ von $(\varepsilon,\eta)^{\nu}$ die Permutation $\gamma^{\nu}$ und die Indices der Coëfficienten $G$ des Productes $(\varepsilon,\eta)^{\nu}(\varepsilon^{\lambda_1},\eta)^{\nu_1}(\varepsilon^{\lambda_2},\eta)^{\nu_2}\dots$ die Permutation $\gamma^{\nu+\lambda_1\nu_1+\lambda_2\nu_2\dots}$}
